\chapter{Projet guidé : la calculatrice (et plus)}

\begin{synthese}[Au programme]
Mener un projet de programmation du début à la fin : du \textbf{cahier des charges}
au programme testé. On produit une \textbf{calculatrice} complète en C (menu, quatre
opérations, gestion de la division par zéro), puis un mini-logiciel de
\textbf{gestion des notes d'une classe} (tableau de structures). C'est la synthèse
de toute l'année : conditions, boucles, fonctions, tableaux et structures réunis.
\end{synthese}

\section{La démarche de projet}

Jusqu'ici, tu as appris des « briques » : variables, \code{if}, boucles, fonctions,
tableaux, structures. Un \textbf{projet} consiste à assembler ces briques pour
construire un logiciel qui rend un service. On ne se jette pas sur le clavier : on
suit une \textbf{démarche}.

\begin{definition}
La \textbf{démarche de projet} est l'enchaînement organisé des étapes qui mènent
d'un besoin exprimé à un programme fonctionnel et fiable. Les cinq étapes sont :
\textbf{cahier des charges}, \textbf{analyse}, \textbf{algorithme}, \textbf{codage},
\textbf{test et débogage}.
\end{definition}

\subsection{Les cinq étapes}

\begin{enumerate}[label=\arabic*)]
  \item \textbf{Cahier des charges} — On écrit, en français, \emph{ce que le
        programme doit faire} : les entrées, les traitements, les sorties, les cas
        particuliers. C'est le contrat. Exemple : « le programme doit additionner
        deux nombres saisis par l'utilisateur et afficher le résultat ».
  \item \textbf{Analyse} — On découpe le problème en sous-problèmes plus simples
        (méthode \emph{descendante}). Chaque sous-problème deviendra souvent une
        \textbf{fonction}.
  \item \textbf{Algorithme} — On écrit la solution en pseudo-code ou par un
        organigramme, sans se soucier du langage. On vérifie la logique « sur le
        papier ».
  \item \textbf{Codage} — On traduit l'algorithme dans le langage choisi (ici, le
        \textbf{C}), en respectant la syntaxe.
  \item \textbf{Test et débogage} — On exécute le programme avec des jeux d'essai
        (valeurs normales \emph{et} cas limites), on repère les erreurs
        (\emph{bugs}) et on les corrige.
\end{enumerate}

\begin{methode}
\textbf{Réussir un projet de programmation.}
(1)~Lis le sujet et \textbf{reformule} le besoin avec tes mots ;
(2)~liste les \textbf{entrées}, le \textbf{traitement} et les \textbf{sorties} ;
(3)~découpe en \textbf{fonctions} (une fonction = une tâche) ;
(4)~code et compile \textbf{morceau par morceau} ;
(5)~teste chaque cas, surtout les \textbf{cas limites} (zéro, valeur négative,
saisie vide).
\end{methode}

\begin{remarque}
Un bon réflexe : compiler souvent. Un programme qui grossit de 200 lignes avant la
première compilation est très difficile à déboguer. Mieux vaut écrire 10 lignes,
compiler, vérifier, puis continuer.
\end{remarque}

\subsection{Repérer entrées, traitement, sorties}

Avant tout projet, on remplit ce petit tableau. Pour notre calculatrice :

\begin{center}\small
\begin{tabular}{@{}lll@{}}
\toprule
\textbf{Entrées} & \textbf{Traitement} & \textbf{Sorties}\\
\midrule
Le choix de l'opération & Choisir la bonne opération & Le résultat du calcul\\
Le premier nombre & Calculer & Un message d'erreur si\\
Le second nombre & (gérer la division par 0) & division par zéro\\
\bottomrule
\end{tabular}
\end{center}

\section{Projet 1 — La calculatrice en C}

\subsection{Cahier des charges}

\begin{definition}
Le \textbf{cahier des charges} décrit précisément le comportement attendu du
logiciel. C'est le document de référence : si le programme le respecte, le travail
est réussi.
\end{definition}

\noindent\textbf{Cahier des charges de la calculatrice :}
\begin{enumerate}[label=\arabic*)]
  \item Le programme affiche un \textbf{menu} : 1) Addition, 2) Soustraction,
        3) Multiplication, 4) Division, 5) Quitter.
  \item L'utilisateur \textbf{choisit} une opération en tapant son numéro.
  \item Pour une opération de calcul, le programme demande \textbf{deux nombres}.
  \item Il \textbf{affiche le résultat} de l'opération.
  \item Pour la division, si le second nombre est \textbf{zéro}, il affiche un
        message d'erreur et ne calcule pas.
  \item Après chaque opération, le programme \textbf{recommence} (réaffiche le
        menu) tant que l'utilisateur n'a pas choisi « Quitter ».
\end{enumerate}

\subsection{Conception (analyse + algorithme)}

On découpe le problème. Le menu doit se répéter : c'est une \textbf{boucle}. Comme
le menu doit s'afficher \emph{au moins une fois}, on choisit une boucle
\code{do\dots while}. À l'intérieur, selon le choix, on exécute une opération
différente : c'est un \code{switch}.

Chaque opération devient une \textbf{fonction} : \code{addition}, \code{soustraction},
\code{multiplication}, \code{division}. Une fonction qui prend deux nombres et
renvoie un résultat est un bon découpage.

\begin{methode}
\textbf{Structure générale (pseudo-code).}
\begin{enumerate}[label=\arabic*)]
  \item Répéter :
  \item \quad afficher le menu ;
  \item \quad lire \code{choix} ;
  \item \quad selon \code{choix} : 1\,$\to$\,addition, 2\,$\to$\,soustraction,
        3\,$\to$\,multiplication, 4\,$\to$\,division, 5\,$\to$\,quitter ;
  \item Tant que \code{choix} $\ne$ 5.
\end{enumerate}
\end{methode}

\begin{attention}
La \textbf{division par zéro} est interdite en mathématiques \emph{et} en
informatique : diviser par 0 provoque un comportement indéfini (souvent un plantage
du programme). Il faut \textbf{toujours} tester le diviseur avant de diviser.
\end{attention}

\subsection{Le code complet commenté}

On utilise le type \code{double} pour accepter les nombres décimaux (ex. 14,5).
Voici le programme complet, prêt à compiler avec \code{gcc calculatrice.c}.

\begin{codec}
/* calculatrice.c -- calculatrice a menu (OnBuch, Tle TI) */
#include <stdio.h>

/* --- Une fonction par operation : chacune renvoie un double --- */
double addition(double a, double b) {
    return a + b;
}

double soustraction(double a, double b) {
    return a - b;
}

double multiplication(double a, double b) {
    return a * b;
}

/* La division renvoie a/b SEULEMENT si b est non nul. */
double division(double a, double b) {
    return a / b;
}

int main(void) {
    int choix;       /* le numero choisi dans le menu     */
    double a, b;     /* les deux nombres a calculer        */

    do {
        /* 1. Afficher le menu */
        printf("\n===== CALCULATRICE OnBuch =====\n");
        printf("1. Addition\n");
        printf("2. Soustraction\n");
        printf("3. Multiplication\n");
        printf("4. Division\n");
        printf("5. Quitter\n");
        printf("Votre choix : ");
        scanf("%d", &choix);

        /* 2. Si on calcule (choix 1 a 4), lire les deux nombres */
        if (choix >= 1 && choix <= 4) {
            printf("Premier nombre  : ");
            scanf("%lf", &a);
            printf("Second nombre   : ");
            scanf("%lf", &b);
        }

        /* 3. Choisir l'operation */
        switch (choix) {
            case 1:
                printf("Resultat : %.2f\n", addition(a, b));
                break;
            case 2:
                printf("Resultat : %.2f\n", soustraction(a, b));
                break;
            case 3:
                printf("Resultat : %.2f\n", multiplication(a, b));
                break;
            case 4:
                /* GESTION DE LA DIVISION PAR ZERO */
                if (b == 0) {
                    printf("Erreur : division par zero impossible !\n");
                } else {
                    printf("Resultat : %.2f\n", division(a, b));
                }
                break;
            case 5:
                printf("Au revoir !\n");
                break;
            default:
                printf("Choix invalide, recommencez.\n");
        }
    } while (choix != 5);   /* 4. Recommencer tant que != 5 */

    return 0;
}
\end{codec}

\fig{Calculatrice complète : menu, fonctions, switch et gestion de la division par zéro.}

\begin{remarque}
On lit un \code{double} avec \code{\%lf} dans \code{scanf} (le « l » est obligatoire),
mais on l'affiche avec \code{\%f} (ou \code{\%.2f} pour deux décimales) dans
\code{printf}. C'est une distinction classique qui piège les débutants.
\end{remarque}

\subsection{Tests}

On vérifie le programme avec des \textbf{jeux d'essai} : des valeurs choisies pour
couvrir le cas normal \emph{et} les cas limites.

\begin{center}\small
\begin{tabular}{@{}llll@{}}
\toprule
\textbf{Choix} & \textbf{a} & \textbf{b} & \textbf{Résultat attendu}\\
\midrule
1 (addition) & 7 & 3 & 10,00\\
2 (soustraction) & 7 & 3 & 4,00\\
3 (multiplication) & 7 & 3 & 21,00\\
4 (division) & 7 & 2 & 3,50\\
4 (division) & 7 & 0 & message d'erreur\\
9 (inexistant) & --- & --- & « Choix invalide »\\
5 (quitter) & --- & --- & fin du programme\\
\bottomrule
\end{tabular}
\end{center}

\begin{methode}
\textbf{Déboguer.} Si un résultat est faux : (1)~vérifie le \textbf{format}
(\code{\%lf} en lecture, \code{\%f} en affichage) ; (2)~ajoute des \code{printf}
temporaires pour afficher les valeurs de \code{a}, \code{b}, \code{choix} ;
(3)~relis l'algorithme : la logique est-elle bonne ?
\end{methode}

\subsection{Améliorations possibles}

Le projet de base fonctionne ; on peut l'enrichir :
\begin{itemize}
  \item ajouter le \textbf{modulo} (reste de la division entière) ;
  \item ajouter la \textbf{puissance} ($a^b$) ;
  \item calculer la \textbf{moyenne} de deux nombres ;
  \item enregistrer l'historique des calculs dans un tableau.
\end{itemize}
Ces extensions sont l'objet des exercices à la fin du chapitre.

\section{Projet 2 — Gestion des notes d'une classe}

Second projet de synthèse : un mini-logiciel pour un professeur du \emph{Lycée
Bilingue de Bafoussam}. Il gère les moyennes d'une classe à l'aide d'un
\textbf{tableau de structures}.

\subsection{Cahier des charges}

\begin{enumerate}[label=\arabic*)]
  \item Chaque élève est décrit par un \textbf{nom} et une \textbf{moyenne}.
  \item Le programme propose un \textbf{menu} : 1) Ajouter un élève, 2) Afficher la
        liste, 3) Moyenne de la classe, 4) Major (meilleur élève), 5) Quitter.
  \item On peut enregistrer jusqu'à 50 élèves.
  \item Le programme recommence tant qu'on n'a pas choisi « Quitter ».
\end{enumerate}

\subsection{Conception}

On définit une \textbf{structure} \code{Eleve} regroupant le nom et la moyenne.
Les élèves sont rangés dans un \textbf{tableau} de \code{Eleve}. Une variable
\code{nb} compte combien d'élèves ont été saisis. Chaque action du menu devient une
\textbf{fonction}.

\begin{syntaxe}
\code{struct Eleve \{ char nom[30]; float moyenne; \};}\quad puis un tableau :
\code{struct Eleve classe[50];}
\end{syntaxe}

\begin{remarque}
On passe le tableau \code{classe} et le nombre \code{nb} aux fonctions. En C, un
tableau passé à une fonction peut être \textbf{modifié} par celle-ci (il est passé
par son adresse). C'est pourquoi \code{ajouter} peut renvoyer le nouveau \code{nb}.
\end{remarque}

\subsection{Le code complet commenté}

\begin{codec}
/* notes.c -- gestion des notes d'une classe (OnBuch, Tle TI) */
#include <stdio.h>
#include <string.h>

#define MAX 50   /* nombre maximal d'eleves */

/* Un eleve : un nom et une moyenne */
struct Eleve {
    char  nom[30];
    float moyenne;
};

/* Ajoute un eleve, renvoie le nouveau nombre d'eleves */
int ajouter(struct Eleve classe[], int nb) {
    if (nb >= MAX) {
        printf("Classe pleine !\n");
        return nb;
    }
    printf("Nom de l'eleve   : ");
    scanf("%s", classe[nb].nom);
    printf("Moyenne sur 20   : ");
    scanf("%f", &classe[nb].moyenne);
    return nb + 1;   /* un eleve de plus */
}

/* Affiche la liste des eleves */
void afficher(struct Eleve classe[], int nb) {
    int i;
    if (nb == 0) {
        printf("Aucun eleve enregistre.\n");
        return;
    }
    printf("\n--- Liste de la classe ---\n");
    for (i = 0; i < nb; i++) {
        printf("%d. %-15s %.2f/20\n", i + 1, classe[i].nom, classe[i].moyenne);
    }
}

/* Calcule et affiche la moyenne de la classe */
void moyenneClasse(struct Eleve classe[], int nb) {
    int i;
    float somme = 0;
    if (nb == 0) {
        printf("Aucun eleve : moyenne impossible.\n");
        return;
    }
    for (i = 0; i < nb; i++) {
        somme = somme + classe[i].moyenne;
    }
    printf("Moyenne de la classe : %.2f/20\n", somme / nb);
}

/* Affiche le major (meilleure moyenne) */
void major(struct Eleve classe[], int nb) {
    int i, indiceMax;
    if (nb == 0) {
        printf("Aucun eleve enregistre.\n");
        return;
    }
    indiceMax = 0;                       /* on suppose le 1er meilleur */
    for (i = 1; i < nb; i++) {
        if (classe[i].moyenne > classe[indiceMax].moyenne) {
            indiceMax = i;               /* on a trouve mieux */
        }
    }
    printf("Major : %s avec %.2f/20\n",
           classe[indiceMax].nom, classe[indiceMax].moyenne);
}

int main(void) {
    struct Eleve classe[MAX];   /* le tableau d'eleves */
    int nb = 0;                 /* nombre d'eleves saisis */
    int choix;

    do {
        printf("\n===== GESTION DES NOTES =====\n");
        printf("1. Ajouter un eleve\n");
        printf("2. Afficher la liste\n");
        printf("3. Moyenne de la classe\n");
        printf("4. Major de la classe\n");
        printf("5. Quitter\n");
        printf("Votre choix : ");
        scanf("%d", &choix);

        switch (choix) {
            case 1: nb = ajouter(classe, nb);   break;
            case 2: afficher(classe, nb);       break;
            case 3: moyenneClasse(classe, nb);  break;
            case 4: major(classe, nb);          break;
            case 5: printf("Au revoir !\n");    break;
            default: printf("Choix invalide.\n");
        }
    } while (choix != 5);

    return 0;
}
\end{codec}

\fig{Gestion des notes : structure Eleve, tableau, et une fonction par action du menu.}

\begin{attention}
\code{scanf("\%s", classe[nb].nom)} lit un mot \textbf{sans espace} : tape
\code{Awa} et non \code{Awa Bello}. Pour gérer les prénoms composés, il faudrait
\code{fgets}, hors programme ici : on s'en tient à un seul mot.
\end{attention}

\begin{remarque}
Dans \code{scanf("\%f", \&classe[nb].moyenne)} on met le \code{\&} (adresse) car on
lit \emph{vers} une variable. Mais pour \code{classe[nb].nom}, qui est un tableau
de caractères, on n'écrit \textbf{pas} de \code{\&} : le nom d'un tableau est déjà
une adresse.
\end{remarque}

\section{Exercices}

\rubrique{Application}

\exo{Étendre la calculatrice : le modulo}\diff{1}
Ajoute au Projet 1 une option « 6. Modulo » qui affiche le reste de la division
entière de \code{a} par \code{b} (opérateur \code{\%} sur des entiers). Que faut-il
prévoir si \code{b} vaut 0 ?

\exo{Étendre la calculatrice : la puissance}\diff{2}
Ajoute une option « 7. Puissance » qui calcule $a^b$ ($b$ entier positif) à l'aide
d'une fonction \code{puissance(double a, int b)} utilisant une boucle (sans
\code{pow}).

\exo{Mention au bulletin}\diff{1}
Dans le Projet 2, écris une fonction \code{mention(float m)} qui affiche la mention
selon la moyenne : $\geq 16$ « Très Bien », $\geq 14$ « Bien », $\geq 12$
« Assez Bien », $\geq 10$ « Passable », sinon « Insuffisant ». Appelle-la dans
\code{afficher}.

\exo{Compter les admis}\diff{2}
Ajoute au Projet 2 une option « Nombre d'admis » : une fonction \code{compterAdmis}
qui renvoie le nombre d'élèves ayant une moyenne $\geq 10$, et affiche le résultat
ainsi que le taux de réussite en pourcentage.

\rubrique{Entraînement}

\exo{Trier les élèves par moyenne}\diff{3}
Écris une fonction \code{trier(struct Eleve classe[], int nb)} qui range les élèves
de la meilleure à la moins bonne moyenne par un \textbf{tri par sélection}. Ajoute
une option de menu « Classement » qui trie puis affiche la liste.

\exo{Historique de la calculatrice}\diff{3}
Modifie la calculatrice pour qu'elle enregistre chaque résultat dans un tableau
\code{double historique[100]}. Ajoute une option « Afficher l'historique » qui
liste tous les calculs effectués depuis le lancement.

\exo{Moyenne pondérée}\diff{2}
Une matière a une note et un coefficient. Définis une structure \code{Matiere} avec
un nom, une note et un coefficient (entier), saisis 4 matières et calcule la moyenne
pondérée $\dfrac{\sum (\text{note}\times\text{coef})}{\sum \text{coef}}$.

\exo{Recherche d'un élève}\diff{2}
Dans le Projet 2, ajoute une option « Rechercher » qui demande un nom et affiche la
moyenne de l'élève correspondant, ou « Élève introuvable » sinon (parcours du
tableau + \code{strcmp}).

\section{Corrigés}

\corrige{de l'exercice 1}
Le modulo n'existe que pour les \textbf{entiers} en C (opérateur \code{\%}). On
ajoute un \code{case 6} et on protège contre le diviseur nul, exactement comme la
division.
\begin{codec}
case 6:
    if ((int)b == 0) {
        printf("Erreur : modulo par zero impossible !\n");
    } else {
        printf("Reste : %d\n", (int)a % (int)b);
    }
    break;
\end{codec}
On pense aussi à ajouter la ligne \code{printf("6. Modulo\\n");} dans le menu et à
modifier le test de lecture en \code{if (choix >= 1 \&\& choix <= 6)}.

\corrige{de l'exercice 2}
La puissance se calcule en multipliant \code{a} par lui-même \code{b} fois.
\begin{codec}
double puissance(double a, int b) {
    double resultat = 1;   /* a^0 vaut 1 */
    int i;
    for (i = 1; i <= b; i++) {
        resultat = resultat * a;
    }
    return resultat;
}
\end{codec}
Dans \code{main}, on ajoute le \code{case 7} : \code{printf("Resultat : \%.2f\\n",
puissance(a, (int)b));}.

\corrige{de l'exercice 3}
Une cascade de \code{if\dots else if} testée de la plus haute mention à la plus
basse.
\begin{codec}
void mention(float m) {
    if      (m >= 16) printf(" -> Tres Bien\n");
    else if (m >= 14) printf(" -> Bien\n");
    else if (m >= 12) printf(" -> Assez Bien\n");
    else if (m >= 10) printf(" -> Passable\n");
    else              printf(" -> Insuffisant\n");
}
\end{codec}
Dans \code{afficher}, après chaque ligne d'élève, on appelle
\code{mention(classe[i].moyenne);}.

\corrige{de l'exercice 4}
On compte les moyennes $\geq 10$, puis on calcule le pourcentage.
\begin{codec}
int compterAdmis(struct Eleve classe[], int nb) {
    int i, admis = 0;
    for (i = 0; i < nb; i++) {
        if (classe[i].moyenne >= 10) {
            admis = admis + 1;
        }
    }
    return admis;
}
\end{codec}
Appel dans le menu :
\begin{codec}
int a = compterAdmis(classe, nb);
printf("Admis : %d / %d\n", a, nb);
if (nb > 0) {
    printf("Taux de reussite : %.1f%%\n", 100.0 * a / nb);
}
\end{codec}

\corrige{de l'exercice 5}
\textbf{Tri par sélection} : à chaque tour, on cherche le maximum dans la partie non
triée et on l'échange avec le premier élément non trié.
\begin{codec}
void trier(struct Eleve classe[], int nb) {
    int i, j, iMax;
    struct Eleve temp;
    for (i = 0; i < nb - 1; i++) {
        iMax = i;
        for (j = i + 1; j < nb; j++) {
            if (classe[j].moyenne > classe[iMax].moyenne) {
                iMax = j;
            }
        }
        /* echange classe[i] et classe[iMax] */
        temp = classe[i];
        classe[i] = classe[iMax];
        classe[iMax] = temp;
    }
}
\end{codec}
L'échange de deux structures se fait simplement avec une variable temporaire
\code{temp} de type \code{struct Eleve}. Dans le menu, on appelle \code{trier} puis
\code{afficher}.

\begin{aretenir}
Un projet se mène par \textbf{étapes} : cahier des charges $\to$ analyse $\to$
algorithme $\to$ codage $\to$ test. On \textbf{découpe} en fonctions (une fonction
= une tâche), on \textbf{compile souvent} et on \textbf{teste les cas limites}
(division par zéro, classe vide). C'est cette méthode, plus que le code lui-même,
qui fait le bon programmeur.
\end{aretenir}
